% !TEX TS-program = xelatex
% !TEX encoding = UTF-8 Unicode

\documentclass[11pt,a4paper]{report}

% --- Packages ---
\usepackage{geometry}
\geometry{top=2.5cm, bottom=2.5cm, left=2.5cm, right=2.5cm}

\usepackage{amsmath,amssymb}
\usepackage{graphicx}
\usepackage{xcolor}
\usepackage{fancyvrb}
\usepackage{tcolorbox}
\usepackage{booktabs}
\usepackage{longtable}
\usepackage{array}
\usepackage{hyperref}

% --- Colors Definition ---
\definecolor{primaryblue}{RGB}{26, 115, 232}
\definecolor{darknavy}{RGB}{13, 37, 63}
\definecolor{codegray}{RGB}{245, 247, 250}
\definecolor{bordergray}{RGB}{220, 224, 230}
\definecolor{accentgreen}{RGB}{46, 125, 50}
\definecolor{accentred}{RGB}{198, 40, 40}

% --- Hyperref Setup ---
\hypersetup{
    colorlinks=true,
    linkcolor=primaryblue,
    filecolor=magenta,
    urlcolor=primaryblue,
    citecolor=accentgreen,
    pdftitle={مستندات جامع فنی ربات تحلیلگر یوتیوب},
    pdfauthor={تیم توسعه هوش مصنوعی}
}

% --- Custom Info Box ---
\newtcolorbox{infobox}[1]{
    colback=codegray,
    colframe=primaryblue,
    fonttitle=\bfseries,
    coltitle=white,
    title=#1,
    arc=3mm,
    boxrule=1pt
}

\newtcolorbox{warningbox}[1]{
    colback=white,
    colframe=accentred,
    fonttitle=\bfseries,
    coltitle=white,
    title=#1,
    arc=3mm,
    boxrule=1pt
}

% --- Persian Configuration (xepersian must be last) ---
\usepackage{xepersian}
\settextfont[Scale=1.0]{Tahoma}
\setlatintextfont[Scale=0.95]{Arial}

\title{\Huge\textbf{مستندات جامع فنی و معماری سامانه}\\
\LARGE\textbf{ربات هوشمند تحلیل و پردازش ویدیوهای یوتیوب}\\
\large\textbf{(YouTube Analyzer Telegram Bot)}}

\author{\textbf{مهندسی سیستم‌های ابری و یادگیری عمیق}\\
مبتنی بر \lr{Cloudflare Workers} و \lr{Google Gemini AI}}

\date{\today}

\begin{document}

\maketitle

\tableofcontents
\newpage

% =========================================================================
\chapter{مقدمه و معماری کلان سامانه}
% =========================================================================

\section{چکیده اجرایی}
سامانه تحلیلگر ویدیوی یوتیوب، یک بستر تمام‌ابری بدون سرور (\lr{Serverless}) بر بستر \lr{Cloudflare Workers} است که با تکیه بر پیام‌رسان تلگرام به عنوان رابط کاربری (\lr{Frontend}) و مدل‌های هوش مصنوعی چندوجهی \lr{Google Gemini}، امکان استخراج، پردازش، خلاصه‌سازی و شرح صحنه‌به‌صحنه ویدیوهای یوتیوب را به زبان فارسی فراهم می‌سازد.

این سیستم با غلبه بر محدودیت‌های رایج پلتفرم تلگرام (مانند سقف ۴۰۹۶ کاراکتری پیام‌ها) و چالش‌های مصرف منابع در مدل‌های زبانی بزرگ (همچون \lr{Front-Loading Bias} و سقف توکن‌های در دقیقه یا \lr{TPM})، راهکاری پایدار، مقیاس‌پذیر و فوق‌سریع ارائه می‌دهد.

\begin{infobox}{معماری سه‌لایه سامانه}
\begin{enumerate}
    \item \textbf{لایه ارتباطی (رابط تلگرام):} دریافت وب‌هوک‌ها، مدیریت تعامل با کاربر و نمایش پیام‌های چندبخشی.
    \item \textbf{لایه پردازش لبه (ورکر کلودفلر):} نرمال‌سازی پیوندها، اعتبارسنجی کاربر، زمان‌بندی کلیدهای \lr{API} و کشینگ داده‌ها.
    \item \textbf{لایه هوش مصنوعی (Google Gemini):} دریافت فریم‌های ویدیو و صوت، پردازش معنایی چندوجهی و تولید متن فارسی منطبق با ساختار زمانی.
\end{enumerate}
\end{infobox}

\section{ویژگی‌ها و دستاوردهای کلیدی سیستم}
\begin{itemize}
    \item \textbf{سه حالت تحلیلی متنوع:}
    \begin{itemize}
        \item \textbf{خلاصه کوتاه و سریع:} استخراج جوهره و پیام کلیدی در ۲ الی ۳ پاراگراف فشرده.
        \item \textbf{تحلیل کامل و جامع:} کالبدشکافی ساختاری همراه با استدلال زنجیره‌ای تفکر (\lr{Chain of Thought}) در کادرهای تاشونده مدرن تلگرام (\lr{Expandable Blockquote}).
        \item \textbf{شرح ویدیو:} روایت داستانی گام‌به‌گام با پیوندهای کلیک‌پذیر زمانی (\lr{Timestamps}) به ثانیه دقیق ویدیو در یوتیوب.
    \end{itemize}
    \item \textbf{تثبیت زمانی دقیق (\lr{Duration Grounding}):} رفع نقص پایان زودهنگام تحلیل و اجبار هوش مصنوعی به پوشش یکنواخت تا انتهای ویدیو.
    \item \textbf{استراتژی کشینگ بهینه (\lr{KV Cache}):} ذخیره‌سازی ۳۰ روزه نتایج و تحویل آنی (کمتر از ۵۰ میلی‌ثانیه) در درخواست‌های مکرر بدون مصرف توکن.
    \item \textbf{مدیریت خطای سهمیه (\lr{Multi-Key Failover}):} چرخش و جایگزینی خودکار بین کلیدهای متعدد در مواجهه با خطای \lr{HTTP 429}.
    \item \textbf{امنیت و کنترل دسترسی (\lr{Whitelist System}):} مکانیزم تایید دستی کاربران جدید توسط مدیر با قابلیت پاسخ بلادرنگ اینلاین.
\end{itemize}

% =========================================================================
\chapter{چرخه حیات پردازش درخواست (Request Lifecycle)}
% =========================================================================

\section{مراحل گام‌به‌گام از دریافت پیوند تا تحویل خروجی}
\begin{enumerate}
    \item \textbf{دریافت وب‌هوک و اعتبارسنجی مقدماتی:} دریافت درخواست \lr{POST} از تلگرام و بررسی لیست سفید.
    \item \textbf{استخراج شناسه ویدیو (\lr{Video ID}):} نرمال‌سازی لینک با \lr{RegEx} و استخراج شناسه ۱۱ رقمی.
    \item \textbf{واکشی فراداده‌های ویدیو (\lr{Metadata}):} استخراج عنوان، سازنده و طول دقیق ویدیو به ثانیه (\lr{lengthSeconds}).
    \item \textbf{ارزیابی کش محلی (\lr{KV Cache Lookup}):} تحویل آنی در کمتر از ۵۰ میلی‌ثانیه در صورت وجود در کش.
    \item \textbf{پیکربندی پرامپت براساس زمان:} محاسبه زمان دقیق به فرمت دقیقه و ثانیه و تزریق به پرامپت.
    \item \textbf{ارسال درخواست چندوجهی به Google Gemini:} ارسال مستقیم آدرس ویدیو در شیء \lr{fileData} با \lr{mimeType: "video/*"}.
    \item \textbf{قالب‌بندی و تبدیل به HTML سازگار با تلگرام:} استخراج تایم‌استمپ‌ها و تبدیل به لینک کلیک‌پذیر.
    \item \textbf{قطعه‌بندی و ارسال چندمرحله‌ای (\lr{Chunking}):} تقسیم متون بالای ۳۹۰۰ کاراکتر بر اساس مرز پاراگراف‌ها.
\end{enumerate}

% =========================================================================
\chapter{ارتباط با Google Gemini API و محاسبات توکن}
% =========================================================================

\section{ساختار درخواست ورودی (Request Payload)}
\begin{latin}
\begin{Verbatim}[frame=single,fontsize=\small,numbers=left]
{
  "contents": [{
    "parts": [
      {
        "fileData": {
          "fileUri": "https://www.youtube.com/watch?v=s1Xxxxxxx",
          "mimeType": "video/*"
        }
      },
      { "text": "سیستم پرامپت و دستورالعمل با زمان ویدیو..." }
    ]
  }],
  "generationConfig": { "temperature": 0.3, "topK": 40, "topP": 0.95 }
}
\end{Verbatim}
\end{latin}

\section{محاسبه ریاضی توکن‌های مصرفی ویدیو و محدودیت TPM}
طبق مستندات رسمی \lr{Google AI Studio}:
\begin{itemize}
    \item \textbf{فریم‌های تصویر (Video Frames):} در هر ثانیه ۱ فریم استخراج می‌شود. هر فریم بین ۱۰۰ الی ۲۶۳ توکن مصرف می‌کند.
    \item \textbf{جریان صوتی (Audio Track):} معادل ۳۲ توکن در ثانیه مصرف می‌نماید.
    \item \textbf{مجموع نرخ مصرف متوسط:} حدود ۲۹۵ توکن در ثانیه (معادل ۱۷٬۷۰۰ توکن در دقیقه).
\end{itemize}

\begin{warningbox}{علت قطع تحلیل در دقیقه ۱۴ الی ۱۵ در پلن رایگان}
یک ویدیوی ۱۴ دقیقه‌ای حدود ۲۴۷٬۸۰۰ توکن ورودی تولید می‌کند که به سقف ۲۵۰٬۰۰۰ توکن در دقیقه (\lr{TPM}) پلن رایگان گوگل می‌رسد و موجب رخداد خطای \lr{HTTP 429} می‌گردد. با تکنیک \lr{Duration Grounding} و \lr{Failover Pool} این چالش مدیریت شده است.
\end{warningbox}

% =========================================================================
\chapter{مهندسی پرامپت و تثبیت زمانی (Duration Grounding)}
% =========================================================================

\section{علت‌شناسی پدیده انحراف به جلو (Front-Loading Bias)}
مدل‌های زبانی در مواجهه با ویدیوهای بلند، تمایل دارند تمام تمرکز و بودجه توکن خروجی خود را به ۱۰ تا ۱۵ دقیقه اول اختصاص دهند.

\section{پیاده‌سازی Duration Grounding}
\begin{latin}
\begin{Verbatim}[frame=single,fontsize=\small]
const durationSec = meta.duration || 0;
const durationFormatted = durationSec > 0 ? formatDuration(durationSec) : "";
const durationMinutes = durationSec > 0 ? Math.round(durationSec / 60) : 0;

// تزریق زمان در پرامپت:
// مدت زمان کامل این اثر: 23:40 (حدود 24 دقیقه) است.
// قانون حیاتی توزیع زمانی: کل ویدیو باید به طور متوازن پوشش داده شود.
// روایت را از [00:00] آغاز کرده و با فواصل منظم تا ثانیه‌های پایانی ادامه بده.
\end{Verbatim}
\end{latin}

% =========================================================================
\chapter{قطعه‌بندی پیام‌ها و تطبیق با محدودیت‌های تلگرام}
% =========================================================================

\section{الگوریتم هوشمند قطعه‌بندی (Split Text Chunks)}
تلگرام اجازه ارسال پیام‌های بلندتر از ۴۰۹۶ کاراکتر را نمی‌دهد. تابع \lr{\texttt{splitTextChunks}} با سقف ۳۹۰۰ کاراکتر متن را با اولویت‌بندی زیر تفکیک می‌کند:
\begin{enumerate}
    \item شکستن بر اساس مرز دو خط تیره (\lr{\texttt{\textbackslash n\textbackslash n}} - پایان پاراگراف).
    \item شکستن بر اساس تک‌خط (\lr{\texttt{\textbackslash n}}).
    \item شکستن در پایان جمله (\lr{\texttt{. }}).
    \item حفظ تمامیت تگ‌های \lr{HTML}.
\end{enumerate}

% =========================================================================
\chapter{پایگاه داده لبه (Cloudflare Workers KV)}
% =========================================================================

\begin{table}[htbp]
\centering
\caption{کلیدهای ذخیره‌سازی در پایگاه داده لبه}
\begin{tabular}{p{5.5cm}cp{5cm}}
\toprule
\textbf{فرمت نام کلید} & \textbf{مدت اعتبار (TTL)} & \textbf{توضیحات و کاربرد} \\
\midrule
\texttt{cache\_\{mode\}\_\{videoId\}} & ۳۰ روز & ذخیره تحلیل آماده جهت تحویل فوری با صفر توکن \\
\texttt{whitelist\_user\_\{id\}} & دائمی & مشخص‌کننده مجاز بودن کاربر در سیستم \\
\texttt{stats\_daily\_\{date\}\_\{key\}} & ۷ روز & شمارنده تعداد مصرف هر کلید Gemini \\
\texttt{logs\_system} & دائمی (مدور) & ۵۰ لاگ اخیر خطاها و وقایع حساس \\
\bottomrule
\end{tabular}
\end{table}

\end{document}
